\documentclass[tikz,border=10pt]{standalone}
\input{../../docs/tikz/dag_preamble.tex}

\begin{document}

% Agent visual validation loop:
% 1. Compile this file to PDF after every substantive DAG edit.
% 2. Inspect the PDF or a rendered PNG, not just the TeX source.
% 3. Increase spacing or reroute arrows until no box, legend, note, edge, or label overlaps.
%
% Intended lifecycle:
% - During paper intake, replace the scaffold below with every named result in the
%   paper and the dependency arrows between them.
% - After intake, treat this as a stable topology: update node styles/status/text
%   as proofs close, but avoid adding boxes or arrows unless the original named
%   result inventory or dependency map was actually incomplete.
%
% Arrow direction rule:
% Draw edges from prerequisite/source node to dependent/target node.  The
% `dag_arrow` and `dag_dashed_arrow` styles use an explicit end-arrow (`->`),
% so the arrowhead should always land on the node that consumes the dependency.
% Do not reverse the coordinate order just to make routing easier; use anchors
% or an orthogonal route through empty space instead.

\begin{tikzpicture}[
  x=1cm,
  y=1cm,
  every node/.append style={outer sep=4pt}
]

\dagPaperMetadata{Iterative Local Voting for Collective Decision-making in Continuous Spaces}{N. Garg, V. Kamble, A. Goel, D. Marn, K. Munagala}{JAIR 64:315--355}{2019}{https://doi.org/10.1613/jair.1.11358}

% --- Legend: README status vocabulary plus formalized node-type styles. ---
\dagPaperLegendRightOfMetadata{
\node[dag_result, dag_template_legend] (legRes) at (0,0) {formalized\\result};
\node[dag_lemma, dag_template_legend] (legLem) at (4.2,0) {formalized\\lemma};
\node[dag_model, dag_template_legend] (legDef) at (8.4,0) {formalized\\definition};
\node[dag_caveat_legend] (legCav) at (12.6,0) {formalized\\with caveat};
\node[dag_partial, dag_template_legend] (legPart) at (0,-2.6) {partially\\formalized};
\node[dag_conditional, dag_template_legend] (legCond) at (4.2,-2.6) {conditional};
\node[dag_scaffold, dag_template_legend] (legScaf) at (8.4,-2.6) {scaffold};
\node[dag_unformalized, dag_template_legend] (legNot) at (12.6,-2.6) {not started /\\not formalized};
}
\daglegend{(legRes)(legLem)(legDef)(legCav)(legPart)(legCond)(legScaf)(legNot)}{Legend}

\begin{dagPaperBody}

\node[dag_model] (C123) at (0,-4.2) {
  \textbf{C1--C3 / Algorithm 1} \\
  Convex bounded space, unique ideals, bounded density, radius $r_0/t$
};
\node[dag_model] (UtilityDefs) at (6.4,-4.2) {
  \textbf{Definitions 1--3} \\
  $L^p$ utilities, weighted Euclidean utilities, decomposable utilities
};
\node[dag_lemma] (FiniteNorms) at (12.8,-4.2) {
  \textbf{Finite-Norm / Recurrence Layer} \\
  $L^p/L^q$ gradients, subgradients, projections, SSGM shape
};

\node[dag_partial] (SmallRegion) at (0,-7.6) {
  \textbf{Lemmas 1, 2, 4} \\
  Null reductions formalized; small-region estimates are boundary inputs
};
\node[dag_unformalized] (SSGM) at (6.4,-7.6) {
  \textbf{Single SSGM / Certificate Boundary} \\
  Source-to-input, C3 interpretation, stochastic convergence
};
\node[dag_lemma] (DualGradient) at (12.8,-7.6) {
  \textbf{Lemma 3} \\
  Holder-dual gradient norm, derivative, subgradient bridge
};

\node[dag_conditional] (Thm1) at (0,-11.0) {
  \textbf{Theorem 1} \\
  $L^p$ utilities, Model A/B, three dual cases converge to SO
};
\node[dag_conditional] (Thm2) at (6.4,-11.0) {
  \textbf{Theorem 2} \\
  Model B finite Holder-dual norms converge to SO
};
\node[dag_conditional] (Prop1) at (12.8,-11.0) {
  \textbf{Proposition 1} \\
  Weighted Euclidean utilities, $L^2$ neighborhoods converge to SO
};

\node[dag_conditional] (Prop2) at (3.2,-14.4) {
  \textbf{Proposition 2} \\
  Decomposable utilities, $L^\infty$ neighborhoods converge to medians
};
\node[dag_partial] (Drift) at (9.6,-14.4) {
  \textbf{Theorem 3 Escape Boundary} \\
  Selected finite-dot concentration plus projected trace facts remain;
  slack is trace-derived residual geometry
};
\node[dag_conditional] (Thm3) at (9.6,-17.8) {
  \textbf{Theorem 3} \\
  Any convergent Model B/$L^2$ trajectory reaches directional equilibrium
};

\draw[dag_arrow] (C123) -- (SmallRegion);
\draw[dag_arrow] (C123) -- (SSGM);
\draw[dag_arrow] (C123) -- (Drift);
\draw[dag_arrow] (UtilityDefs) -- (SmallRegion);
\draw[dag_arrow] (UtilityDefs) -- (DualGradient);
\draw[dag_arrow] (UtilityDefs) -- (Prop1);
\draw[dag_arrow] (FiniteNorms) -- (SmallRegion);
\draw[dag_arrow] (FiniteNorms) -- (DualGradient);
\draw[dag_dashed_arrow] (SmallRegion) -- (Thm1);
\draw[dag_dashed_arrow] (SSGM) -- (Thm1);
\draw[dag_dashed_arrow] (SSGM) -- (Thm2);
\draw[dag_dashed_arrow] (DualGradient) -- (Thm2);
\draw[dag_dashed_arrow] (SmallRegion) -- (Prop1);
\draw[dag_dashed_arrow] (SSGM) -- (Prop1);
\draw[dag_dashed_arrow] (Thm1) -- (Prop2);
\draw[dag_dashed_arrow] (Drift) -- (Thm3);
\end{dagPaperBody}
\end{tikzpicture}
\end{document}
